\chapter{需求分析}

\section{用户角色}

项目通过 \code{user_account.role} 区分 \code{admin}、\code{student}、\code{teacher} 三类角色。登录认证由 \code{services/auth_service.py} 实现，页面访问控制由 \code{pages/_guards.py} 中的 \code{require_role()} 完成。三类用户需求如下：

\begin{itemize}
  \item 学生：查看当前学期可选课程，提交选课，查看已选课程，按规则退课，查看成绩单和绩点。
  \item 教师：查看本人负责的开课班次，录入或修改学生成绩，查看成绩分布和成绩修改日志。
  \item 管理员：维护学期、课程、开课班次、学生、教师，查看系统统计和所有班次成绩。
\end{itemize}

\section{功能需求}

系统功能需求和真实代码对应关系见表 \ref{tab:function-requirements}。

{\footnotesize
\setlength{\tabcolsep}{3pt}
\renewcommand{\arraystretch}{1.16}
\begin{longtable}{p{2.7cm}p{6.7cm}p{3.9cm}}
\caption{功能需求与项目文件对应关系}\label{tab:function-requirements}\\
\toprule
功能 & 当前实现 & 主要文件 \\
\midrule
\endfirsthead
\toprule
功能 & 当前实现 & 主要文件 \\
\midrule
\endhead
登录与会话 & 用户名密码登录，Cookie 持久会话，按角色进入不同页面 & \code{services/auth_service.py}、\code{utils/auth_cookie.py}、\code{opengauss_setup/sql/add_user_session.sql} \\
课程查询 & 学生、教师、管理员均可通过服务层查询课程与开课班次 & \code{services/course_service.py} \\
选课 & 检查重复、容量、学期状态、选课窗口、先修课、时间冲突后写入 \code{enrollment} & \code{services/selection_service.py} \\
退课 & 只能退本人已选且未录入成绩的课程，并检查学期和时间窗口 & \code{services/selection_service.py} \\
课程与教学班管理 & 管理员维护课程、学期、开课班次和上课时间片 & \code{services/course_service.py}、\code{pages/course_manage.py} \\
成绩管理 & 教师只能维护本人班次，管理员可维护全部班次，并记录成绩日志 & \code{services/score_service.py}、\code{pages/score_manage.py} \\
基础资料维护 & 管理学生、教师资料并同步创建账号 & \code{services/student_service.py}、\code{services/teacher_service.py} \\
\bottomrule
\end{longtable}
}

\section{数据需求}

系统需要保存的数据分为五类。第一类是组织结构数据，包括 \code{department}、\code{major}、\code{classroom}。第二类是用户身份数据，包括 \code{user_account}、\code{user_session}、\code{student}、\code{teacher}、\code{admin_profile}。第三类是教学安排数据，包括 \code{semester}、\code{course}、\code{course_offering}、\code{course_schedule} 和 \code{course_prerequisite}。第四类是选课成绩数据，即 \code{enrollment}。第五类是审计数据，即 \code{score_change_log}。

\section{业务规则}

当前已经实现的业务规则包括：

\begin{enumerate}
  \item 重复选课限制：\code{select_course_tx()} 会锁定并检查同一学生、同一教学班的已有选课记录，\code{enrollment} 表上的 \code{UNIQUE(student_id, offering_id)} 作为数据库层兜底。
  \item 容量限制：\code{trg_enrollment_guard} 在锁定 \code{course_offering} 行后统计 \code{status='selected'} 的人数，并与 \code{max_capacity} 比较。
  \item 行级锁保护：选课函数和选课触发器均使用 \code{SELECT ... FOR UPDATE} 锁定目标 \code{course_offering} 行。
  \item 时间冲突检查：\code{trg_enrollment_guard} 通过 \code{course_schedule} 的 \code{weekday}、\code{start_section}、\code{end_section} 判断同学期区间重叠。
  \item 先修课限制：\code{trg_enrollment_guard} 通过 \code{course_prerequisite} 和历史 \code{completed} 成绩判断先修课是否达到 60 分。
  \item 退课规则：只能退本人 \code{selected} 状态且未录入成绩的记录；数据库触发器禁止已完成或已有成绩的记录退课。
  \item 成绩权限：教师只能维护自己负责的班次，管理员可以维护所有班次。
  \item 成绩范围：数据库 CHECK 与服务层均限制成绩为 0 到 100。
  \item 成绩审计：\code{trg_score_change_log} 在 \code{final_score} 变化后自动写入 \code{score_change_log}，操作者由成绩更新事务传入。
  \item 角色档案一致性：学生、教师、管理员档案由触发器检查其 \code{user_id} 对应账户角色是否匹配。
  \item 教学班与教室容量一致性：\code{trg_offering_capacity_check} 和 \code{trg_classroom_capacity_update_check} 维护 \code{max_capacity} 与教室容量之间的跨表约束。
\end{enumerate}

仍需进一步完善的规则包括：同一学生同一学期选择同一课程不同教学班的限制、数据库级角色授权、备份恢复流程以及多连接并发压测。

\section{数据字典草案}

\begin{longtable}{p{3.4cm}p{3.8cm}p{7cm}}
\caption{核心数据字典草案}\label{tab:data-dictionary}\\
\toprule
数据项 & 来源表 & 含义 \\
\midrule
\endfirsthead
\toprule
数据项 & 来源表 & 含义 \\
\midrule
\endhead
\code{student_id} & \code{student} & 学生业务编号，学生主键。 \\
\code{teacher_id} & \code{teacher} & 教师业务编号，教师主键。 \\
\code{course_id} & \code{course} & 课程编号，表示课程定义。 \\
\code{offering_id} & \code{course_offering} & 教学班代理主键，表示某学期某教师开设的具体课次。 \\
\code{schedule_id} & \code{course_schedule} & 上课时间片主键。 \\
\code{enrollment_id} & \code{enrollment} & 选课记录主键。 \\
\code{final_score} & \code{enrollment} & 学生在某教学班的最终成绩。 \\
\code{token_hash} & \code{user_session} & 持久登录 token 的哈希值。 \\
\bottomrule
\end{longtable}
